% ============================================================================
% APPENDIX GB: LIT-BECKER - Gary Becker Research Integration
% ============================================================================
% Category: LIT (Literature Integration)
% Version: 1.0
% Last Updated: 2026-01-17
% Dependencies: K (LIT-FEHR), L (LIT-KAHNEMAN), BBB (CORE-WHERE)
% Language: english
% ============================================================================

% ----------------------------------------------------------------------------
% HEADER BLOCK
% ----------------------------------------------------------------------------
% Code: GB
% Category: LIT (Literature Integration)
% Title: Gary Becker Research Integration
% Author: EBF Framework
% Status: Complete
% Priority: High
% Evidence-Base: 30 papers (see data/paper-sources.yaml)
% ----------------------------------------------------------------------------

\chapter{LIT-BECKER: Gary Becker Research Integration}
\label{app:lit-becker}


\begin{tcolorbox}[colback=blue!5!white, colframe=blue!75!black,
    title=Quick Reference: Appendix UNMAPPED_GB]

\textbf{Appendix:} GB (LIT)

\textbf{Core Question:} What are the key findings in this domain?

\textbf{Key Concepts:}
\begin{itemize}[nosep]
\item Framework integration
\item Parameter validation
\item Cross-references
\end{itemize}

\textbf{Cross-References:} See related appendices below
\end{tcolorbox}

% ============================================================================
% CROSS-REFERENCE MAP
% ============================================================================
\begin{tcolorbox}[colback=blue!5!white,colframe=blue!75!black,title=Cross-Reference Map]
\textbf{Outgoing References:}
\begin{itemize}
    \item CORE-WHO (AAA): Welfare hierarchy and household decision-making
    \item CORE-WHAT (C): Utility dimensions beyond consumption
    \item CORE-HOW (B): Complementarity in human capital and addiction
    \item CORE-WHEN (V): Time allocation and lifecycle dynamics
    \item CORE-WHERE (BBB): Empirical parameters for education returns
    \item CORE-AWARE (AU): Information in discrimination models
    \item CORE-READY (AV): Investment willingness in human capital
    \item CORE-STAGE (AW): Lifecycle stages and habit formation
    \item LIT-FEHR (K): Social preferences vs rational choice
    \item LIT-KAHNEMAN (L): Bounded rationality comparison
    \item LIT-THALER (M): Mental accounting and self-control
    \item METHOD-LLMMC (AN): Parameter estimation for returns
    \item DOMAIN-HEALTH (AA): Health capital and addiction
    \item DOMAIN-EDUCATION (AB): Education investment models
\end{itemize}

\textbf{Incoming References:}
\begin{itemize}
    \item All DOMAIN appendices: Rational choice foundations
    \item FORMAL-AXIOMS (A): Utility maximization axioms
    \item METHOD-ESTIMATION (AL): Structural parameter estimation
\end{itemize}
\end{tcolorbox}

% ============================================================================
% CHAPTER LINKAGE
% ============================================================================
\begin{tcolorbox}[colback=green!5!white,colframe=green!75!black,title=Chapter Linkage]
\textbf{Primary Chapter:} Chapter 3 (Rational Choice Foundations)

\textbf{Supporting Chapters:}
\begin{itemize}
    \item Chapter 5: Time preferences and discounting
    \item Chapter 7: Family and household economics
    \item Chapter 9: Human capital accumulation
    \item Chapter 14: Health behaviors and addiction
\end{itemize}
\end{tcolorbox}

% ============================================================================
% ABSTRACT AND QUICK REFERENCE
% ============================================================================
\begin{abstract}
This appendix integrates Gary S. Becker's foundational contributions to economics into the EBF framework. Becker (Nobel Prize 1992) revolutionized economics by extending rational choice analysis to domains previously considered outside economics: human capital investment, discrimination, crime, family behavior, addiction, and time allocation. His work provides the rational choice benchmark against which behavioral economics identifies systematic deviations. The EBF framework incorporates Becker's insights through six axioms (UNMAPPED_GB-1 to UNMAPPED_GB-6) that formalize human capital complementarity, household production, rational addiction, and the economics of discrimination. Key parameters include education returns ($r \approx 8$--$12\%$), adjacent complementarity in addiction ($\partial^2 U/\partial c_t \partial S > 0$), and household time allocation efficiency. Becker's framework remains essential for understanding when behavioral interventions are necessary versus when standard incentives suffice.
\end{abstract}

\begin{tcolorbox}[colback=gray!5!white,colframe=gray!75!black,title=Quick Reference]
\textbf{Core Contributions:}
\begin{itemize}
    \item Human capital theory: Education as investment with measurable returns
    \item Economics of discrimination: Taste-based vs statistical discrimination
    \item Crime and punishment: Rational criminal choice under uncertainty
    \item Family economics: Household production and marriage markets
    \item Rational addiction: Intertemporal complementarity in consumption
    \item Time allocation: Household production function approach
\end{itemize}

\textbf{Key Parameters:}
\begin{itemize}
    \item Education returns: $r \approx 8$--$12\%$ per year of schooling
    \item Discrimination coefficient: $d > 0$ raises minority hiring costs
    \item Crime elasticity: Probability effect $> $ severity effect
    \item Addiction complementarity: $\gamma_{c,S} > 0$ (adjacent complementarity)
    \item Time value: Shadow price $w^* = \partial H/\partial t_h \cdot p_Z$
\end{itemize}

\textbf{Evidence Tier:} Mixed (Tier 1-3)
\begin{itemize}
    \item Human capital returns: Tier 1 (IV identification via schooling laws)
    \item Discrimination: Tier 2 (audit studies, natural experiments)
    \item Rational addiction: Tier 2-3 (time series, quasi-experiments)
    \item Crime deterrence: Tier 1-2 (RCTs, discontinuity designs)
\end{itemize}
\end{tcolorbox}

% ============================================================================
% FUNDAMENTAL QUESTION
% ============================================================================
% -----------------------------------------------------------------------------
% APPENDIX SCOPE BOX
% -----------------------------------------------------------------------------

\begin{tcolorbox}[
    colback=orange!5!white,
    colframe=orange!75!black,
    title={\textbf{Appendix Scope: Ziel / In-Scope / Out-of-Scope / Constraints}},
    fonttitle=\bfseries
]

\textbf{Ziel (Objective):}
\begin{itemize}[nosep]
    \item How can the rational choice framework be extended to explain behaviors traditionally considered outside the domain of economics---and when do systematic deviations from this framework require behavioral interventions?
\end{itemize}

\textbf{In-Scope:}
\begin{itemize}[nosep]
    \item Fundamental Question
    \item Key Results and Evidence
    \item EBF Integration
    \item Worked Example: Education Investment with Present Bias
    \item Critical Assessment
\end{itemize}

\textbf{Out-of-Scope (delegiert an andere Appendices):}
\begin{itemize}[nosep]
    \item REF-GLOSSARY $\rightarrow$ Appendix G
    \item REF-GLOSSARY $\rightarrow$ Appendix G
    \item CORE-WHERE $\rightarrow$ Appendix UNMAPPED_BBB
\end{itemize}

\textbf{Constraints (Anwendungsgrenzen):}
\begin{itemize}[nosep]
    \item [TODO: Define when this appendix does NOT apply]
    \item [TODO: Define required prerequisites]
    \item [TODO: Define known limitations]
\end{itemize}

\textbf{Lieferobjekte (Deliverables):}
\begin{itemize}[nosep]
    \item 6 Table(s)
    \item 16 Equation(s)
    \item 6 Axiom(s)
    \item Worked Example(s)
\end{itemize}

\end{tcolorbox}


\section{Fundamental Question}
\label{sec:gb-fundamental}

\begin{tcolorbox}[colback=red!5!white,colframe=red!75!black,title=Fundamental Question]
\textbf{How can the rational choice framework be extended to explain behaviors traditionally considered outside the domain of economics---and when do systematic deviations from this framework require behavioral interventions?}
\end{tcolorbox}

Gary Becker's research program asked a revolutionary question: Can economic analysis---based on maximizing behavior, stable preferences, and market equilibrium---explain human behavior in domains like education, family, crime, and addiction? His affirmative answer expanded economics' scope dramatically while providing the benchmark against which behavioral economics measures deviations.

For the EBF framework, Becker's contributions are essential because they:
\begin{enumerate}
    \item Establish the rational choice baseline that behavioral interventions must outperform
    \item Provide structural parameters (education returns, time values) used throughout EBF
    \item Identify when price/incentive mechanisms suffice versus when behavioral tools are needed
    \item Formalize complementarity structures that EBF extends to context-dependent preferences
\end{enumerate}

% ============================================================================
% THEORETICAL FOUNDATIONS
% ============================================================================
\section{Theoretical Foundations}
\label{sec:gb-theory}

\subsection{Human Capital Theory}
\label{subsec:gb-human-capital}

Becker's human capital theory \citep{becker1964human} treats education and training as investment decisions with measurable returns:

\begin{equation}
\max_{s} \sum_{t=0}^{T} \frac{w(s,t) - c(s,t)}{(1+r)^t}
\label{eq:gb-human-capital}
\end{equation}

where $s$ is years of schooling, $w(s,t)$ is wage at time $t$ given schooling $s$, $c(s,t)$ is direct and opportunity costs, and $r$ is the discount rate.

The Mincer equation operationalizes this:
\begin{equation}
\ln w = \alpha + \beta_s \cdot s + \beta_1 \cdot x + \beta_2 \cdot x^2 + \epsilon
\label{eq:gb-mincer}
\end{equation}

where $\beta_s \approx 0.08$--$0.12$ represents the return to an additional year of schooling.

\begin{axiom}[Human Capital Investment]
\label{ax:gb-1}
\textbf{UNMAPPED_GB-1:} Individuals invest in human capital until the marginal return equals the marginal cost:
\begin{equation}
\frac{\partial PV(w)}{\partial s} = MC(s)
\end{equation}
Investment is increasing in ability, decreasing in discount rates, and responsive to expected returns. Deviations from optimal investment indicate either market failures (credit constraints) or behavioral frictions (present bias, information gaps).
\end{axiom}

\textbf{EBF Integration:} Axiom UNMAPPED_GB-1 connects to CORE-READY (AV) through investment willingness and to CORE-AWARE (AU) through information about returns. When individuals underinvest relative to UNMAPPED_GB-1 predictions despite adequate information and credit access, behavioral interventions targeting present bias (commitment devices) or salience (return information) may be warranted.

\subsection{Economics of Discrimination}
\label{subsec:gb-discrimination}

Becker's discrimination model \citep{becker1957economics} introduces a ``taste for discrimination'' coefficient $d$:

\begin{equation}
\text{Effective wage cost} = w_m(1 + d)
\label{eq:gb-discrimination}
\end{equation}

where $w_m$ is the minority wage and $d > 0$ represents the discriminator's willingness to pay to avoid minority workers.

\begin{axiom}[Discrimination and Market Competition]
\label{ax:gb-2}
\textbf{UNMAPPED_GB-2:} Taste-based discrimination is costly to discriminators and should diminish under competitive pressure:
\begin{equation}
\frac{\partial d^*}{\partial \text{competition}} < 0
\end{equation}
Statistical discrimination based on group characteristics persists even under competition when information is costly. The distinction between taste-based and statistical discrimination determines appropriate policy responses.
\end{axiom}

\textbf{EBF Integration:} Axiom UNMAPPED_GB-2 connects to CORE-AWARE (AU) through information asymmetries in statistical discrimination. Behavioral interventions (blind review, structured interviews) address taste-based discrimination; information interventions address statistical discrimination.

\subsection{Crime and Punishment}
\label{subsec:gb-crime}

Becker's economic approach to crime \citep{becker1968crime} models criminal choice as:

\begin{equation}
EU_{\text{crime}} = (1-p) \cdot U(Y + G) + p \cdot U(Y - F) - U(Y)
\label{eq:gb-crime}
\end{equation}

where $p$ is probability of punishment, $G$ is gain from crime, $F$ is punishment severity, and $Y$ is legitimate income.

\begin{axiom}[Rational Criminal Choice]
\label{ax:gb-3}
\textbf{UNMAPPED_GB-3:} Criminal behavior responds to incentives---both probability and severity of punishment:
\begin{equation}
\frac{\partial \text{Crime}}{\partial p} < 0, \quad \frac{\partial \text{Crime}}{\partial F} < 0
\end{equation}
Empirically, the probability elasticity exceeds the severity elasticity: $|\epsilon_p| > |\epsilon_F|$. This suggests police presence and detection are more effective than harsher sentences.
\end{axiom}

\textbf{EBF Integration:} Axiom UNMAPPED_GB-3 connects to CORE-WHEN (V) through salience of punishment probability. Behavioral insights suggest probability effects are larger partly because immediate, certain consequences are more salient than delayed, uncertain ones---consistent with present bias (CORE-HOW).

\subsection{Family Economics}
\label{subsec:gb-family}

Becker's household production model \citep{becker1965theory, becker1981treatise} treats the family as a production unit:

\begin{equation}
\max U(Z_1, Z_2, \ldots, Z_n) \quad \text{s.t.} \quad Z_i = f_i(x_i, t_i)
\label{eq:gb-household}
\end{equation}

where $Z_i$ are ``commodities'' (meals, child quality, leisure activities) produced using market goods $x_i$ and time $t_i$.

\begin{axiom}[Household Production Efficiency]
\label{ax:gb-4}
\textbf{UNMAPPED_GB-4:} Households allocate time and goods to maximize production efficiency:
\begin{equation}
\frac{\partial Z/\partial t_i}{\partial Z/\partial x_i} = \frac{w}{p}
\end{equation}
where $w$ is the wage rate (opportunity cost of time) and $p$ is the price of market goods. Comparative advantage determines household specialization.
\end{axiom}

\textbf{EBF Integration:} Axiom UNMAPPED_GB-4 connects to CORE-WHO (AAA) through household-level welfare aggregation. The EBF framework extends this by incorporating behavioral frictions in household bargaining, mental accounting of household budgets, and present bias in household savings decisions.

\subsection{Rational Addiction}
\label{subsec:gb-addiction}

Becker and Murphy's rational addiction model \citep{becker1988theory} introduces adjacent complementarity:

\begin{equation}
\max \sum_{t=0}^{\infty} \beta^t U(c_t, S_t) \quad \text{s.t.} \quad S_{t+1} = (1-\delta)S_t + c_t
\label{eq:gb-addiction}
\end{equation}

where $c_t$ is addictive consumption, $S_t$ is the addiction stock, and $\delta$ is the depreciation rate.

\begin{axiom}[Adjacent Complementarity in Addiction]
\label{ax:gb-5}
\textbf{UNMAPPED_GB-5:} Addictive goods exhibit adjacent complementarity---past consumption increases marginal utility of current consumption:
\begin{equation}
\frac{\partial^2 U}{\partial c_t \partial S_t} > 0 \quad \text{(reinforcement)}
\end{equation}
Combined with $\partial U/\partial S < 0$ (tolerance/withdrawal), this generates the characteristic addiction dynamics: escalating consumption, difficulty quitting, and response to anticipated future prices.
\end{axiom}

\textbf{EBF Integration:} Axiom UNMAPPED_GB-5 is central to EBF's treatment of habit formation and connects directly to CORE-HOW (B) complementarity structures. The key extension: behavioral economics adds present bias ($\beta < 1$), which explains why many addicts want to quit but fail---they are ``rationally addicted'' only under exponential discounting.

\subsection{Time Allocation}
\label{subsec:gb-time}

Becker's time allocation model \citep{becker1965theory} recognizes time as the ultimate scarce resource:

\begin{equation}
T = t_w + \sum_i t_i \quad \text{(time constraint)}
\label{eq:gb-time}
\end{equation}

\begin{equation}
\text{Full income} = wT + V = \sum_i (p_i x_i + w t_i)
\label{eq:gb-full-income}
\end{equation}

where $w$ is the wage, $T$ is total time, $t_w$ is work time, and $V$ is non-labor income.

\begin{axiom}[Time as Binding Constraint]
\label{ax:gb-6}
\textbf{UNMAPPED_GB-6:} As wages rise, time-intensive activities become relatively more expensive, shifting consumption toward goods-intensive activities:
\begin{equation}
\frac{\partial (t_i/x_i)}{\partial w} < 0
\end{equation}
This explains why high-income individuals substitute market goods for time (convenience foods, domestic help) and why time scarcity has increased despite rising incomes.
\end{axiom}

\textbf{EBF Integration:} Axiom UNMAPPED_GB-6 connects to CORE-WHEN (V) through temporal context effects. Behavioral economics extends this by noting that time scarcity itself affects decision quality---bandwidth constraints under time pressure lead to suboptimal choices (Mullainathan \& Shafir).

% ============================================================================
% KEY RESULTS AND EVIDENCE
% ============================================================================
\section{Key Results and Evidence}
\label{sec:gb-results}

\subsection{Human Capital Returns}

\begin{table}[htbp]
\centering
\caption{Returns to Education: Selected Estimates}
\label{tab:gb-returns}
\begin{tabular}{llcc}
\toprule
\textbf{Study} & \textbf{Method} & \textbf{Return (\%)} & \textbf{Evidence Tier} \\
\midrule
Card (1999) & IV (distance to college) & 9.0--13.0 & Tier 1 \\
Angrist \& Krueger (1991) & IV (quarter of birth) & 7.0--8.0 & Tier 1 \\
Oreopoulos (2006) & RD (compulsory schooling) & 10.0--14.0 & Tier 1 \\
Psacharopoulos (2018) & Meta-analysis & 8.8 (global avg) & Tier 2 \\
Mincer (1974) & OLS baseline & 7.0--10.0 & Tier 3 \\
\bottomrule
\end{tabular}
\end{table}

The robust finding: returns to education are approximately 8--12\% per year, with causal estimates often exceeding OLS estimates due to ability bias and measurement error.

\subsection{Discrimination Evidence}

\begin{table}[htbp]
\centering
\caption{Discrimination: Audit Study Evidence}
\label{tab:gb-discrimination-evidence}
\begin{tabular}{llcc}
\toprule
\textbf{Study} & \textbf{Context} & \textbf{Discrimination Gap} & \textbf{Evidence Tier} \\
\midrule
Bertrand \& Mullainathan (2004) & Resume callbacks (US) & 50\% fewer callbacks & Tier 1 \\
Neumark et al. (1996) & Restaurant hiring & 40\% gap (high-price) & Tier 1 \\
Oreopoulos (2011) & Resume callbacks (Canada) & 40\% fewer callbacks & Tier 1 \\
Carlsson \& Rooth (2007) & Resume callbacks (Sweden) & 30\% gap (Arab names) & Tier 1 \\
\bottomrule
\end{tabular}
\end{table}

Audit studies consistently find discrimination, though distinguishing taste-based from statistical discrimination remains challenging.

\subsection{Crime Deterrence}

\begin{table}[htbp]
\centering
\caption{Crime Deterrence: Elasticity Estimates}
\label{tab:gb-crime-evidence}
\begin{tabular}{llcc}
\toprule
\textbf{Study} & \textbf{Intervention} & \textbf{Elasticity} & \textbf{Evidence Tier} \\
\midrule
Levitt (1997) & Police (IV: elections) & $-0.9$ to $-1.0$ & Tier 1 \\
Di Tella \& Schargrodsky (2004) & Police (terror attack) & $-0.75$ & Tier 1 \\
Drago et al. (2009) & Sentence length (Italy) & $-0.05$ to $-0.15$ & Tier 1 \\
Chalfin \& McCrary (2018) & Police meta-analysis & $-0.4$ to $-0.6$ & Tier 2 \\
\bottomrule
\end{tabular}
\end{table}

Consistent with Becker's model, probability of apprehension has larger effects than sentence severity.

\subsection{Rational Addiction Tests}

\begin{table}[htbp]
\centering
\caption{Rational Addiction: Empirical Tests}
\label{tab:gb-addiction-evidence}
\begin{tabular}{llcc}
\toprule
\textbf{Study} & \textbf{Finding} & \textbf{Support for RA} & \textbf{Evidence Tier} \\
\midrule
Becker et al. (1994) & Cigarette price response & Yes (future prices matter) & Tier 2 \\
Gruber \& K\"{o}szegi (2001) & Tax anticipation & Mixed (present bias needed) & Tier 2 \\
Bernheim \& Rangel (2004) & Cue-triggered consumption & No (visceral factors) & Tier 2 \\
Orphanides \& Zervos (1995) & Learning about addiction & Partial (uncertainty) & Tier 3 \\
\bottomrule
\end{tabular}
\end{table}

Evidence for rational addiction is mixed: forward-looking behavior exists but present bias and cue-triggered consumption require behavioral extensions.

% ============================================================================
% EBF INTEGRATION
% ============================================================================
\section{EBF Integration}
\label{sec:gb-integration}

\subsection{CORE Framework Mapping}

\begin{table}[htbp]
\centering
\caption{Becker's Contributions Mapped to 10C Framework}
\label{tab:gb-9c-mapping}
\begin{tabular}{lll}
\toprule
\textbf{CORE} & \textbf{Becker Contribution} & \textbf{Behavioral Extension} \\
\midrule
WHO (AAA) & Household as decision unit & Bargaining frictions, fairness \\
WHAT (C) & Extended utility (Z-goods) & Reference dependence in commodities \\
HOW (B) & Adjacent complementarity & Present bias in addiction \\
WHEN (V) & Time allocation model & Bandwidth constraints \\
WHERE (BBB) & Structural estimation & Behavioral parameter elicitation \\
AWARE (AU) & Information in discrimination & Attention and salience \\
READY (AV) & Investment willingness & Implementation intentions \\
STAGE (AW) & Lifecycle human capital & Behavioral change journey \\
\bottomrule
\end{tabular}
\end{table}

\subsection{When Rational Choice Suffices vs. Behavioral Intervention Needed}

Becker's framework helps identify when standard incentives work:

\begin{tcolorbox}[colback=green!5!white,colframe=green!75!black,title=Rational Choice Sufficient]
\begin{itemize}
    \item High stakes, repeated decisions (UNMAPPED_GB-1: human capital)
    \item Competitive markets with learning opportunities
    \item Clear feedback on outcomes
    \item Long time horizons with commitment mechanisms available
\end{itemize}
\end{tcolorbox}

\begin{tcolorbox}[colback=red!5!white,colframe=red!75!black,title=Behavioral Intervention Needed]
\begin{itemize}
    \item Infrequent decisions with delayed feedback
    \item Decisions under time pressure or cognitive load
    \item Self-control problems (addiction, savings)
    \item Complex information environments
    \item Decisions affecting future selves
\end{itemize}
\end{tcolorbox}

\subsection{Parameter Integration with BBB}

Key parameters from Becker's work for CORE-WHERE (BBB):

\begin{table}[htbp]
\centering
\caption{Becker Parameters in EBF Framework}
\label{tab:gb-parameters}
\begin{tabular}{llcc}
\toprule
\textbf{Parameter} & \textbf{Context} & \textbf{Value} & \textbf{Source} \\
\midrule
$r_{\text{education}}$ & Returns to schooling & 0.08--0.12 & Card (1999) \\
$\epsilon_{\text{police}}$ & Crime-police elasticity & $-0.5$ to $-1.0$ & Meta-analysis \\
$\gamma_{c,S}$ & Addiction complementarity & $> 0$ & Becker \& Murphy \\
$\delta_{\text{health}}$ & Health capital depreciation & 0.02--0.04/year & Grossman (1972) \\
$w^*$ & Shadow value of time & $0.3$--$0.5 \times w$ & Household surveys \\
\bottomrule
\end{tabular}
\end{table}

% ============================================================================
% WORKED EXAMPLE
% ============================================================================
\section{Worked Example: Education Investment with Present Bias}
\label{sec:gb-example}

\begin{tcolorbox}[colback=yellow!5!white,colframe=yellow!75!black,title=Worked Example: College Enrollment Decision]
\textbf{Scenario:} Maria (age 18) is deciding whether to attend college. She faces:
\begin{itemize}
    \item Tuition cost: \$15,000/year for 4 years
    \item Forgone earnings: \$25,000/year
    \item Expected wage premium: 60\% higher lifetime earnings
    \item Discount rate: $r = 5\%$
\end{itemize}

\textbf{Becker Model (UNMAPPED_GB-1):}
\begin{align}
NPV &= \sum_{t=23}^{65} \frac{0.6 \times w_{\text{HS}}}{(1.05)^{t-18}} - \sum_{t=1}^{4} \frac{15,000 + 25,000}{(1.05)^t} \\
&\approx \$500,000 - \$145,000 = \$355,000 > 0
\end{align}

\textbf{Prediction:} Maria should attend college---NPV is strongly positive.

\textbf{Behavioral Extension (Present Bias):}
With quasi-hyperbolic discounting ($\beta = 0.7$, $\delta = 0.95$):
\begin{equation}
NPV_{\text{present}} = \beta \times \$500,000 - \$145,000 = \$350,000 - \$145,000 = \$205,000
\end{equation}

Still positive, but the immediate costs loom larger. If Maria also faces:
\begin{itemize}
    \item Liquidity constraints (no access to loans)
    \item Uncertainty about completion (30\% dropout rate)
    \item First-generation student (no family guidance)
\end{itemize}

\textbf{EBF Intervention Design:}
\begin{enumerate}
    \item \textbf{Address credit constraints:} Income-contingent loans (rational response)
    \item \textbf{Address present bias:} Commitment devices, automatic enrollment
    \item \textbf{Address information gaps:} Simplified FAFSA, return calculators
    \item \textbf{Address identity:} First-gen student programs, role models
\end{enumerate}

\textbf{Expected Effect:} Combined interventions increase enrollment by 15--25 percentage points among marginal students (Bettinger et al., 2012; Hoxby \& Turner, 2015).
\end{tcolorbox}

% ============================================================================
% CRITICAL ASSESSMENT
% ============================================================================
\section{Critical Assessment}
\label{sec:gb-critical}

\subsection{Strengths of Becker's Framework}

\begin{enumerate}
    \item \textbf{Unifying power:} Single framework explains diverse behaviors
    \item \textbf{Testable predictions:} Clear empirical implications
    \item \textbf{Policy relevance:} Direct implications for incentive design
    \item \textbf{Parameter estimation:} Structural approach enables counterfactuals
    \item \textbf{Benchmark function:} Identifies when deviations require explanation
\end{enumerate}

\subsection{Limitations and Behavioral Extensions}

\begin{enumerate}
    \item \textbf{Stable preferences assumption:} Evidence for reference dependence, framing effects
    \item \textbf{Exponential discounting:} Present bias documented across domains
    \item \textbf{Full information:} Bounded attention, salience effects matter
    \item \textbf{Computational capacity:} Complex optimization unrealistic for many decisions
    \item \textbf{Social influences:} Peer effects beyond rational information transmission
\end{enumerate}

\subsection{Synthesis: Becker + Behavioral Economics}

The EBF framework synthesizes Becker's insights with behavioral economics:

\begin{equation}
U_{\text{EBF}} = U_{\text{Becker}}(Z; \Theta) + \Psi(\text{context}) + \epsilon_{\text{behavioral}}
\label{eq:gb-synthesis}
\end{equation}

where $U_{\text{Becker}}$ provides the rational baseline, $\Psi$ captures context effects, and $\epsilon_{\text{behavioral}}$ represents systematic deviations (present bias, social preferences, limited attention).

% ============================================================================
% SUMMARY
% ============================================================================

% === AUTO-GENERATED ENTRIES (2026-02-22) ===

%%%%%%%%%%%%%%%%%%%%%%%%%%%%%%%%%%%%%%%%%%%%%%%%%%%%%%%%%%%%%%%%%%%%%%%%%%%%%%%
\subsubsection{1962: Investment in Human Capital: A Theoretical Analysis}
%%%%%%%%%%%%%%%%%%%%%%%%%%%%%%%%%%%%%%%%%%%%%%%%%%%%%%%%%%%%%%%%%%%%%%%%%%%%%%%

\paragraph{Citation.}
Becker, Gary S (1962). ``Investment in Human Capital: A Theoretical Analysis.'' \textit{Journal of Political Economy}, 70, 9--49.

\paragraph{10C Integration.}
\begin{itemize}[nosep]
  \item \textbf{Domain}: [To be specified]
  \item \textbf{use\_for}: LIT-BECKER
\end{itemize}


%%%%%%%%%%%%%%%%%%%%%%%%%%%%%%%%%%%%%%%%%%%%%%%%%%%%%%%%%%%%%%%%%%%%%%%%%%%%%%%
\subsubsection{1977: De Gustibus Non Est Disputandum}
%%%%%%%%%%%%%%%%%%%%%%%%%%%%%%%%%%%%%%%%%%%%%%%%%%%%%%%%%%%%%%%%%%%%%%%%%%%%%%%

\paragraph{Citation.}
Stigler, George J and Becker, Gary S (1977). ``De Gustibus Non Est Disputandum.'' \textit{American Economic Review}, 67, 76--90.

\paragraph{Core Finding.}
Foundational paper establishing stable preferences assumption and household production approach

\paragraph{10C Integration.}
\begin{itemize}[nosep]
  \item \textbf{Domain}: [To be specified]
  \item \textbf{use\_for}: CORE-WHAT, CORE-HOW, CORE-WHO, CORE-WHEN, CORE-AWARE, METHOD-DESIGN, LIT-STIGLER, LIT-BECKER
\end{itemize}


%%%%%%%%%%%%%%%%%%%%%%%%%%%%%%%%%%%%%%%%%%%%%%%%%%%%%%%%%%%%%%%%%%%%%%%%%%%%%%%
\subsubsection{1992: Habits, Addictions, and Traditions}
%%%%%%%%%%%%%%%%%%%%%%%%%%%%%%%%%%%%%%%%%%%%%%%%%%%%%%%%%%%%%%%%%%%%%%%%%%%%%%%

\paragraph{Citation.}
Becker, Gary S. (1992). ``Habits, Addictions, and Traditions.'' \textit{Kyklos}, 45, 327--345.

\paragraph{10C Integration.}
\begin{itemize}[nosep]
  \item \textbf{Domain}: [To be specified]
  \item \textbf{use\_for}: LIT-BECKER
\end{itemize}


%%%%%%%%%%%%%%%%%%%%%%%%%%%%%%%%%%%%%%%%%%%%%%%%%%%%%%%%%%%%%%%%%%%%%%%%%%%%%%%
\subsubsection{1985: Altruism, Egoism, and Genetic Fitness: Economics and Sociobi...}
%%%%%%%%%%%%%%%%%%%%%%%%%%%%%%%%%%%%%%%%%%%%%%%%%%%%%%%%%%%%%%%%%%%%%%%%%%%%%%%

\paragraph{Citation.}
Becker, Gary S. (1985). ``Altruism, Egoism, and Genetic Fitness: Economics and Sociobiology.'' \textit{Journal of Economic Literature}, 14.

\paragraph{10C Integration.}
\begin{itemize}[nosep]
  \item \textbf{Domain}: [To be specified]
  \item \textbf{use\_for}: LIT-BECKER
\end{itemize}


%%%%%%%%%%%%%%%%%%%%%%%%%%%%%%%%%%%%%%%%%%%%%%%%%%%%%%%%%%%%%%%%%%%%%%%%%%%%%%%
\subsubsection{1985: Public Policies, Pressure Groups, and Dead Weight Costs}
%%%%%%%%%%%%%%%%%%%%%%%%%%%%%%%%%%%%%%%%%%%%%%%%%%%%%%%%%%%%%%%%%%%%%%%%%%%%%%%

\paragraph{Citation.}
Becker, Gary S. (1985). ``Public Policies, Pressure Groups, and Dead Weight Costs.'' \textit{Journal of Public Economics}, 28.

\paragraph{10C Integration.}
\begin{itemize}[nosep]
  \item \textbf{Domain}: [To be specified]
  \item \textbf{use\_for}: LIT-BECKER
\end{itemize}


%%%%%%%%%%%%%%%%%%%%%%%%%%%%%%%%%%%%%%%%%%%%%%%%%%%%%%%%%%%%%%%%%%%%%%%%%%%%%%%
\subsubsection{1991: An Analysis of Cigarette Addiction}
%%%%%%%%%%%%%%%%%%%%%%%%%%%%%%%%%%%%%%%%%%%%%%%%%%%%%%%%%%%%%%%%%%%%%%%%%%%%%%%

\paragraph{Citation.}
Becker, Gary S. and Murphy, Kevin M. (1991). ``An Analysis of Cigarette Addiction.'' \textit{Working Paper}.

\paragraph{10C Integration.}
\begin{itemize}[nosep]
  \item \textbf{Domain}: [To be specified]
  \item \textbf{use\_for}: LIT-BECKER
\end{itemize}


%%%%%%%%%%%%%%%%%%%%%%%%%%%%%%%%%%%%%%%%%%%%%%%%%%%%%%%%%%%%%%%%%%%%%%%%%%%%%%%
\subsubsection{1991: A Note on Restaurant Pricing and Other Examples of Social In...}
%%%%%%%%%%%%%%%%%%%%%%%%%%%%%%%%%%%%%%%%%%%%%%%%%%%%%%%%%%%%%%%%%%%%%%%%%%%%%%%

\paragraph{Citation.}
Becker, Gary S. and Murphy, Kevin M. (1991). ``A Note on Restaurant Pricing and Other Examples of Social Influences on Price.'' \textit{Journal of Political Economy}, 99.

\paragraph{10C Integration.}
\begin{itemize}[nosep]
  \item \textbf{Domain}: [To be specified]
  \item \textbf{use\_for}: LIT-BECKER
\end{itemize}


%%%%%%%%%%%%%%%%%%%%%%%%%%%%%%%%%%%%%%%%%%%%%%%%%%%%%%%%%%%%%%%%%%%%%%%%%%%%%%%
\subsubsection{1988: A Reformulation of the Economic Theory of Fertility}
%%%%%%%%%%%%%%%%%%%%%%%%%%%%%%%%%%%%%%%%%%%%%%%%%%%%%%%%%%%%%%%%%%%%%%%%%%%%%%%

\paragraph{Citation.}
Becker, Gary S. and Barro, Robert J. (1988). ``A Reformulation of the Economic Theory of Fertility.'' \textit{Quarterly Journal of Economics}, 103.

\paragraph{10C Integration.}
\begin{itemize}[nosep]
  \item \textbf{Domain}: [To be specified]
  \item \textbf{use\_for}: LIT-BECKER
\end{itemize}


%%%%%%%%%%%%%%%%%%%%%%%%%%%%%%%%%%%%%%%%%%%%%%%%%%%%%%%%%%%%%%%%%%%%%%%%%%%%%%%
\subsubsection{1993: Nobel Lecture: The Economic Way of Looking at Behavior}
%%%%%%%%%%%%%%%%%%%%%%%%%%%%%%%%%%%%%%%%%%%%%%%%%%%%%%%%%%%%%%%%%%%%%%%%%%%%%%%

\paragraph{Citation.}
Becker, Gary S. (1993). ``Nobel Lecture: The Economic Way of Looking at Behavior.'' \textit{Journal of Political Economy}, 101.

\paragraph{10C Integration.}
\begin{itemize}[nosep]
  \item \textbf{Domain}: [To be specified]
  \item \textbf{use\_for}: LIT-BECKER
\end{itemize}


%%%%%%%%%%%%%%%%%%%%%%%%%%%%%%%%%%%%%%%%%%%%%%%%%%%%%%%%%%%%%%%%%%%%%%%%%%%%%%%
\subsubsection{1994: Human Capital: A Theoretical and Empirical Analysis (3rd Edi...}
%%%%%%%%%%%%%%%%%%%%%%%%%%%%%%%%%%%%%%%%%%%%%%%%%%%%%%%%%%%%%%%%%%%%%%%%%%%%%%%

\paragraph{Citation.}
Becker, Gary S. (1994). \textit{Human Capital: A Theoretical and Empirical Analysis (3rd Edition)}. Publisher.

\paragraph{10C Integration.}
\begin{itemize}[nosep]
  \item \textbf{Domain}: [To be specified]
  \item \textbf{use\_for}: LIT-BECKER
\end{itemize}


%%%%%%%%%%%%%%%%%%%%%%%%%%%%%%%%%%%%%%%%%%%%%%%%%%%%%%%%%%%%%%%%%%%%%%%%%%%%%%%
\subsubsection{1996: Accounting for Tastes}
%%%%%%%%%%%%%%%%%%%%%%%%%%%%%%%%%%%%%%%%%%%%%%%%%%%%%%%%%%%%%%%%%%%%%%%%%%%%%%%

\paragraph{Citation.}
Becker, Gary S. (1996). \textit{Accounting for Tastes}. Publisher.

\paragraph{10C Integration.}
\begin{itemize}[nosep]
  \item \textbf{Domain}: [To be specified]
  \item \textbf{use\_for}: LIT-BECKER
\end{itemize}


%%%%%%%%%%%%%%%%%%%%%%%%%%%%%%%%%%%%%%%%%%%%%%%%%%%%%%%%%%%%%%%%%%%%%%%%%%%%%%%
\subsubsection{1997: The Endogenous Determination of Time Preference}
%%%%%%%%%%%%%%%%%%%%%%%%%%%%%%%%%%%%%%%%%%%%%%%%%%%%%%%%%%%%%%%%%%%%%%%%%%%%%%%

\paragraph{Citation.}
Becker, Gary S. and Mulligan, Casey B. (1997). ``The Endogenous Determination of Time Preference.'' \textit{Quarterly Journal of Economics}, 112.

\paragraph{10C Integration.}
\begin{itemize}[nosep]
  \item \textbf{Domain}: [To be specified]
  \item \textbf{use\_for}: LIT-BECKER
\end{itemize}


%%%%%%%%%%%%%%%%%%%%%%%%%%%%%%%%%%%%%%%%%%%%%%%%%%%%%%%%%%%%%%%%%%%%%%%%%%%%%%%
\subsubsection{2000: Social Economics: Market Behavior in a Social Environment}
%%%%%%%%%%%%%%%%%%%%%%%%%%%%%%%%%%%%%%%%%%%%%%%%%%%%%%%%%%%%%%%%%%%%%%%%%%%%%%%

\paragraph{Citation.}
Becker, Gary S. and Murphy, Kevin M. (2000). \textit{Social Economics: Market Behavior in a Social Environment}. Publisher.

\paragraph{10C Integration.}
\begin{itemize}[nosep]
  \item \textbf{Domain}: [To be specified]
  \item \textbf{use\_for}: LIT-BECKER
\end{itemize}


%%%%%%%%%%%%%%%%%%%%%%%%%%%%%%%%%%%%%%%%%%%%%%%%%%%%%%%%%%%%%%%%%%%%%%%%%%%%%%%
\subsubsection{2005: The Quantity and Quality of Life and the Evolution of World ...}
%%%%%%%%%%%%%%%%%%%%%%%%%%%%%%%%%%%%%%%%%%%%%%%%%%%%%%%%%%%%%%%%%%%%%%%%%%%%%%%

\paragraph{Citation.}
Becker, Gary S. and Philipson, Tomas J. and Soares, Rodrigo R. (2005). ``The Quantity and Quality of Life and the Evolution of World Inequality.'' \textit{American Economic Review}, 95.

\paragraph{10C Integration.}
\begin{itemize}[nosep]
  \item \textbf{Domain}: [To be specified]
  \item \textbf{use\_for}: LIT-BECKER
\end{itemize}


%%%%%%%%%%%%%%%%%%%%%%%%%%%%%%%%%%%%%%%%%%%%%%%%%%%%%%%%%%%%%%%%%%%%%%%%%%%%%%%
\subsubsection{2005: The Equilibrium Distribution of Income and the Market for St...}
%%%%%%%%%%%%%%%%%%%%%%%%%%%%%%%%%%%%%%%%%%%%%%%%%%%%%%%%%%%%%%%%%%%%%%%%%%%%%%%

\paragraph{Citation.}
Becker, Gary S. and Murphy, Kevin M. and Werning, Ivan (2005). ``The Equilibrium Distribution of Income and the Market for Status.'' \textit{Journal of Political Economy}, 113.

\paragraph{10C Integration.}
\begin{itemize}[nosep]
  \item \textbf{Domain}: [To be specified]
  \item \textbf{use\_for}: LIT-BECKER
\end{itemize}


%%%%%%%%%%%%%%%%%%%%%%%%%%%%%%%%%%%%%%%%%%%%%%%%%%%%%%%%%%%%%%%%%%%%%%%%%%%%%%%
\subsubsection{2026: {The Likelihood of Persistently Low Global Fertility}}
%%%%%%%%%%%%%%%%%%%%%%%%%%%%%%%%%%%%%%%%%%%%%%%%%%%%%%%%%%%%%%%%%%%%%%%%%%%%%%%

\paragraph{Citation.}
Geruso, Michael and Spears, Dean (2026). ``{The Likelihood of Persistently Low Global Fertility}.'' \textit{Journal of Economic Perspectives}, 40, 3--26.

\paragraph{Core Finding.}
JEP Winter 2026 symposium paper. Full text available in data/paper-texts/PAP-geruso2026likelihood.md.

\paragraph{Key Parameters.}
\texttt{TFR_global_2025=2.06, TFR_1950=4.85, fertility_persistence_tau}

\paragraph{10C Integration.}
\begin{itemize}[nosep]
  \item \textbf{Domain}: [To be specified]
  \item \textbf{use\_for}: LIT-O, DOMAIN-DEMOGRAPHY, CORE-WHAT, CORE-STAGE, LIT-BECKER
\end{itemize}


%%%%%%%%%%%%%%%%%%%%%%%%%%%%%%%%%%%%%%%%%%%%%%%%%%%%%%%%%%%%%%%%%%%%%%%%%%%%%%%
\subsubsection{2026: {Family Institutions and the Global Fertility Transition}}
%%%%%%%%%%%%%%%%%%%%%%%%%%%%%%%%%%%%%%%%%%%%%%%%%%%%%%%%%%%%%%%%%%%%%%%%%%%%%%%

\paragraph{Citation.}
Gobbi, Paula E. and Hannusch, Anne and Rossi, Pauline (2026). ``{Family Institutions and the Global Fertility Transition}.'' \textit{Journal of Economic Perspectives}, 40, 47--70.

\paragraph{Core Finding.}
JEP Winter 2026 symposium paper. Full text available in data/paper-texts/PAP-gobbi2026family.md.

\paragraph{Key Parameters.}
\texttt{alpha_inheritance, gamma_gender_role, beta_institution, rho_econ_conditional}

\paragraph{10C Integration.}
\begin{itemize}[nosep]
  \item \textbf{Domain}: [To be specified]
  \item \textbf{use\_for}: LIT-O, DOMAIN-FAMILY, CORE-WHAT, CORE-WHEN, LIT-BECKER
\end{itemize}


%%%%%%%%%%%%%%%%%%%%%%%%%%%%%%%%%%%%%%%%%%%%%%%%%%%%%%%%%%%%%%%%%%%%%%%%%%%%%%%
\subsubsection{2026: {Global Labor Mobility between Shrinking and Growing Labor F...}
%%%%%%%%%%%%%%%%%%%%%%%%%%%%%%%%%%%%%%%%%%%%%%%%%%%%%%%%%%%%%%%%%%%%%%%%%%%%%%%

\paragraph{Citation.}
Pritchett, Lant (2026). ``{Global Labor Mobility between Shrinking and Growing Labor Forces}.'' \textit{Journal of Economic Perspectives}, 40, 71--92.

\paragraph{Core Finding.}
JEP Winter 2026 symposium paper. Full text available in data/paper-texts/PAP-pritchett2026global.md.

\paragraph{Key Parameters.}
\texttt{TFR_rich_below_1.5, TFR_africa_above_3.5, dependency_ratio_shift, mobility_gains_trillions}

\paragraph{10C Integration.}
\begin{itemize}[nosep]
  \item \textbf{Domain}: [To be specified]
  \item \textbf{use\_for}: LIT-O, DOMAIN-MIGRATION, CORE-WHEN, CORE-HIERARCHY, LIT-BECKER
\end{itemize}


%%%%%%%%%%%%%%%%%%%%%%%%%%%%%%%%%%%%%%%%%%%%%%%%%%%%%%%%%%%%%%%%%%%%%%%%%%%%%%%
\subsubsection{2026: {How Much Would Continued Low Fertility Affect the US Standa...}
%%%%%%%%%%%%%%%%%%%%%%%%%%%%%%%%%%%%%%%%%%%%%%%%%%%%%%%%%%%%%%%%%%%%%%%%%%%%%%%

\paragraph{Citation.}
Weil, David N. (2026). ``{How Much Would Continued Low Fertility Affect the US Standard of Living?}.'' \textit{Journal of Economic Perspectives}, 40, 27--46.

\paragraph{Core Finding.}
JEP Winter 2026 symposium paper. Full text available in data/paper-texts/PAP-weil2026continued.md.

\paragraph{Key Parameters.}
\texttt{TFR_US_2023=1.62, consumption_impact_TFR1.0=-8.7pct, transition_benefit_decades=4}

\paragraph{10C Integration.}
\begin{itemize}[nosep]
  \item \textbf{Domain}: [To be specified]
  \item \textbf{use\_for}: LIT-O, DOMAIN-DEMOGRAPHY, CORE-WHAT, CORE-STAGE, LIT-BECKER
\end{itemize}

% === END AUTO-GENERATED ENTRIES ===
\section{Summary}
\label{sec:gb-summary}

\begin{tcolorbox}[colback=blue!5!white,colframe=blue!75!black,title=Key Takeaways]
\begin{enumerate}
    \item \textbf{UNMAPPED_GB-1 (Human Capital):} Education investment responds to returns; deviations indicate market or behavioral failures
    \item \textbf{UNMAPPED_GB-2 (Discrimination):} Taste-based vs statistical discrimination require different interventions
    \item \textbf{UNMAPPED_GB-3 (Crime):} Probability of punishment more effective than severity
    \item \textbf{UNMAPPED_GB-4 (Household):} Families optimize production; behavioral frictions in bargaining and budgeting
    \item \textbf{UNMAPPED_GB-5 (Addiction):} Adjacent complementarity explains habit formation; present bias explains failure to quit
    \item \textbf{UNMAPPED_GB-6 (Time):} Rising wages increase time scarcity; bandwidth constraints affect decision quality
\end{enumerate}

\textbf{Integration Principle:} Becker's framework provides the rational choice benchmark; EBF identifies systematic deviations and designs behaviorally-informed interventions when standard incentives fail.
\end{tcolorbox}

% ============================================================================
% GLOSSARY
% ============================================================================
\section{Glossary}
\label{sec:gb-glossary}

Key terms defined in Appendix G (REF-GLOSSARY):

\begin{itemize}
    \item \textbf{Human capital:} Skills, knowledge, and health that affect productivity
    \item \textbf{Adjacent complementarity:} Past consumption increases current marginal utility
    \item \textbf{Taste-based discrimination:} Willingness to pay to avoid certain groups
    \item \textbf{Statistical discrimination:} Using group characteristics as information proxy
    \item \textbf{Household production function:} Technology converting time and goods to commodities
    \item \textbf{Full income:} Value of total time plus non-labor income
    \item \textbf{Shadow price of time:} Implicit value of time in household production
\end{itemize}

% ============================================================================
% CRITICAL FOUNDATIONS
% ============================================================================
\section{Critical Foundations}
\label{sec:gb-foundations}

\begin{enumerate}
    \item Becker, G. S. (1957). \textit{The Economics of Discrimination}. University of Chicago Press.
    \item Becker, G. S. (1964). \textit{Human Capital}. Columbia University Press.
    \item Becker, G. S. (1965). A Theory of the Allocation of Time. \textit{Economic Journal}, 75(299), 493--517.
    \item Becker, G. S. (1968). Crime and Punishment: An Economic Approach. \textit{Journal of Political Economy}, 76(2), 169--217.
    \item Becker, G. S. (1976). \textit{The Economic Approach to Human Behavior}. University of Chicago Press.
    \item Becker, G. S. (1981). \textit{A Treatise on the Family}. Harvard University Press.
    \item Becker, G. S., \& Murphy, K. M. (1988). A Theory of Rational Addiction. \textit{Journal of Political Economy}, 96(4), 675--700.
    \item Becker, G. S., Grossman, M., \& Murphy, K. M. (1994). An Empirical Analysis of Cigarette Addiction. \textit{American Economic Review}, 84(3), 396--418.
    \item Mincer, J. (1974). \textit{Schooling, Experience, and Earnings}. Columbia University Press.
    \item Grossman, M. (1972). On the Concept of Health Capital. \textit{Journal of Political Economy}, 80(2), 223--255.
\end{enumerate}

% ============================================================================
% OPEN ISSUES
% ============================================================================
\section{Open Issues}
\label{sec:gb-open}

\begin{enumerate}
    \item \textbf{Preference stability:} How to reconcile Becker's stable preferences with documented preference reversals?
    \item \textbf{Addiction welfare:} Should we respect revealed preferences of addicts or paternalistic interventions?
    \item \textbf{Discrimination persistence:} Why does discrimination persist despite competitive pressures?
    \item \textbf{Human capital externalities:} How large are social returns beyond private returns?
    \item \textbf{Household bargaining:} How to integrate cooperative and non-cooperative models?
\end{enumerate}

% ============================================================================
% REFERENCES SECTION
% ============================================================================
\section{References}
\label{sec:gb-references}

\nocite{bcm_master}

This appendix integrates 30 papers from Gary Becker's research corpus. Full references are maintained in the master bibliography (\texttt{bibliography/bcm\_master.bib}).

\textbf{Core Becker Works (30 papers):}
\begin{itemize}
    \item \textbf{Human Capital:} Becker (1964, 1967, 1975, 1993)
    \item \textbf{Discrimination:} Becker (1957, 1971)
    \item \textbf{Crime:} Becker (1968), Becker \& Landes (1974)
    \item \textbf{Family:} Becker (1973, 1974, 1981, 1985, 1991)
    \item \textbf{Time Allocation:} Becker (1965), Michael \& Becker (1973)
    \item \textbf{Addiction:} Becker \& Murphy (1988), Becker et al. (1991, 1994)
    \item \textbf{Social Interactions:} Becker (1974), Becker \& Murphy (2000)
    \item \textbf{Economic Approach:} Becker (1976, 1996)
\end{itemize}

\textbf{Key Empirical Tests:}
\begin{itemize}
    \item Card (1999): IV estimates of education returns
    \item Angrist \& Krueger (1991): Quarter-of-birth instruments
    \item Bertrand \& Mullainathan (2004): Resume audit studies
    \item Levitt (1997): Police and crime elasticities
    \item Gruber \& K\"{o}szegi (2001): Present bias in addiction
\end{itemize}

See Appendix G (REF-GLOSSARY) for term definitions and Appendix UNMAPPED_BBB (CORE-WHERE) for parameter values.

\end{document}
